% !TeX root = main.tex
% Judul BAB I (Pusat)
\begin{center}
	\textbf{BAB I} \\
	\textbf{PENDAHULUAN}
\end{center}
\addcontentsline{toc}{section}{BAB I PENDAHULUAN}

\vspace{0.5cm}

% Menggunakan command \hurufsection dari main.tex Anda
\hurufsection{Latar Belakang}{
	Perkembangan teknologi informasi dan komunikasi (TIK) telah mengubah paradigma pengelolaan informasi di seluruh dunia. Perpustakaan, sebagai lembaga pengelola informasi, dituntut untuk senantiasa beradaptasi dengan perubahan tersebut agar tetap relevan dalam melayani kebutuhan pemustaka. Era digital tidak hanya mengubah bentuk koleksi dari cetak menjadi elektronik, tetapi juga mengubah cara masyarakat dalam mengakses dan memanfaatkan data.

	Di Indonesia, transformasi perpustakaan menuju sistem digital menjadi krusial untuk mengatasi hambatan jarak dan waktu. Menurut Putu Laxman Pendit, perpustakaan digital bukan sekadar kumpulan koleksi elektronik, melainkan sebuah sistem yang menyediakan layanan informasi yang terintegrasi dan berkelanjutan.\footcite{pendit2009} Hal ini menuntut adanya kesiapan infrastruktur serta sumber daya manusia yang kompeten di bidang teknologi informasi.

	Namun, dalam praktiknya, masih banyak perpustakaan yang menghadapi kendala dalam mengimplementasikan sistem otomasi dan digitalisasi secara penuh. Kurangnya pemahaman mengenai tata kelola data dan aspek teknis teknologi informasi seringkali menjadi penghambat utama.\footcite{saleh2014} Oleh karena itu, penting untuk membahas lebih dalam mengenai penerapan teknologi informasi di perpustakaan agar tantangan-tantangan tersebut dapat diatasi dengan solusi yang tepat. Berdasarkan fenomena tersebut, makalah ini disusun untuk mengeksplorasi peran teknologi informasi dalam meningkatkan kualitas layanan perpustakaan di masa depan.
}

\hurufsection{Rumusan Masalah}{
	Berdasarkan latar belakang di atas, maka rumusan masalah dalam makalah ini adalah:
	\begin{enumerate}[label=\arabic*., leftmargin=1.5cm, nosep]
		\item Apa yang dimaksud dengan teknologi informasi dalam pendidikan?
		\item Bagaimana dampak penggunaan teknologi terhadap motivasi belajar siswa?
	\end{enumerate}
}

\hurufsection{Tujuan Penulisan}{
	Adapun tujuan dari penulisan makalah ini adalah:
	\begin{enumerate}[label=\arabic*., leftmargin=1.5cm, nosep]
		\item Untuk mengetahui definisi teknologi informasi dalam konteks pendidikan.
		\item Untuk menganalisis pengaruh teknologi terhadap semangat belajar di sekolah.
	\end{enumerate}
}